\documentclass[11pt,a4paper]{article}

\usepackage{iftex}
\ifPDFTeX
  \usepackage[utf8]{inputenc}
  \usepackage[T1]{fontenc}
  \usepackage{lmodern}
\else
  \usepackage{fontspec}
  \setmainfont{Liberation Serif}
  \setsansfont{Liberation Sans}
  \setmonofont{Liberation Mono}
\fi
\usepackage[czech]{babel}
\usepackage{amsmath,amssymb,amsthm,mathtools}
\usepackage[margin=25mm]{geometry}
\usepackage{microtype}
\usepackage[hidelinks]{hyperref}

\theoremstyle{definition}
\newtheorem{definition}{Definice}
\theoremstyle{remark}
\newtheorem{remark}{Poznámka}

\title{\textbf{Rozšíření Eulerova prostoru}\\[4pt]
\large na horizontu účtu}
\author{Jan Mikulík}
\date{19. listopadu 2026}

\begin{document}
\maketitle

\begin{abstract}
Toto je model, ne důkaz.
Eulerův prostor je jednotková kružnice.
Horizont je hranice, za kterou vnitřní účet přestává účtovat.
Rozšíření přidá sterilní směr ortogonálním součtem.
Kdo mění osud uvnitř, platí.
Kdo je jen slyšen zvenčí, vnitřní účet neodečte, a prostor stavů se zvětší.
Z nuly vznikne jednotka.
Sterilní neutrino je jméno tohoto směru, ne jeho existence.
\end{abstract}

\section{Eulerův prostor}

\begin{definition}[Eulerův prostor]
\[
  \mathcal E=\{e^{i\theta}:\theta\in\mathbb R\}\subset\mathbb C.
\]
Zobrazení \(U(\theta)=e^{i\theta}\) splňuje \(|U(\theta)|=1\).
Na kružnici tedy zachovává délku.
\end{definition}

Vnitřní účet na parametru je \(c(\theta)=\theta\).
Posun zevnitř,
\[
  \theta\mapsto\theta+\delta,\qquad \delta\neq 0,
\]
odečte platbu \(|\delta|\).
To je změna osudu uvnitř: akce je v záznamu a účet klesne vůči výchozí hodnotě o tuto platbu.

\begin{definition}[Vznik z nuly]
\[
  e^{i\cdot 0}=1.
\]
Jednotka je tu.
Reálný účet v nule zůstává \(c(0)=0\).
\end{definition}

Z ničeho, měřeno vnitřním účtem, je tu nový bod kružnice.
Samotná formule dodavatele nesestrojuje.
Říká jen, jak takový vznik v tomto jazyce vypadá.

\section{Horizont účtu}

\begin{definition}[Horizont]
Horizont \(\mathcal H\) je hranice oblasti, na které je \(c\) definováno.
Na \(\mathcal H\) a vně něj vnitřní účet neúčtuje.
\end{definition}

\begin{remark}
Není to věta o horizontu černé díry.
Je to jméno hranice účtu.
Fyzikální horizont událostí by byl modelem této hranice teprve po samostatném ztotožnění.
To zde není.
\end{remark}

\section{Rozšíření}

Aktivní prostor je \(\mathcal E_{\mathrm a}\).
Sterilní směr je přímka \(\mathcal E_{\mathrm s}\), která v \(\mathcal E_{\mathrm a}\) neleží.
Rozšířený prostor je ortogonální součet
\[
  \mathcal E^{\sharp}=\mathcal E_{\mathrm a}\oplus\mathcal E_{\mathrm s}.
\]

\begin{definition}[Unitární rozšíření]
Na \(\mathcal E^{\sharp}\) je evoluce unitární,
\[
  U^{\sharp}(\theta)
  =
  \begin{pmatrix}
    e^{i\theta} & 0\\
    0 & 1
  \end{pmatrix},
\]
když se sterilní směr netočí vnitřním parametrem.
Obecněji je přípustná každá unitární mapa na součtu.
Podmínka modelu je jiná: částečný účet na \(\mathcal E_{\mathrm a}\) sterilní složku neodečte.
\end{definition}

\begin{definition}[Vnější slyšení]
Přání \(w\) vybere \(s\in\mathcal E_{\mathrm s}\).
Slyšení bez platby je krok
\[
  x\mapsto x\oplus s
\]
takový, že
\[
  c(x\oplus s)\big|_{\mathcal E_{\mathrm a}}=c(x).
\]
\end{definition}

Prostor stavů se zvětší,
\[
  X_{n+1}=X_n\cup\{s_{n+1}\},
\]
a integrál \(c\) přes původní aktivní množinu se nezmění.
Koza se nažere: nový stav je v záznamu.
Tráva zůstane: vnitřní účet neklesl.

Ortogonální součet není provázání.
Provázání chce netriviální vektor v tenzorovém součinu dvou soustav.
Stav, který leží celé v \(\mathcal E_{\mathrm s}\), není vnitřním účtem vázán k \(\mathcal E_{\mathrm a}\).
Superpozice \(\alpha a+\beta s\) v jednom rozšířeném prostoru je pořád jeden systém, ne pár.
Model tedy říká „neprovázaný a unitární“ v tomto přesném smyslu: evoluce na součtu je unitární a sterilní složka nemá vnitřní účet.

\section{Sterilní neutrino jako čtení}

Aktivní neutrino platí do účtu slabého náboje.
Jeho oscilace je vnitřní posun: chuť se mění a účet aktivního bloku to vidí.

Sterilní neutrino, pokud je ve spektru, nemá slabý náboj.
Do účtu \(c\) slabé interakce neplatí.
Míchání s aktivními může být unitární na větším prostoru chutí.
Na aktivním bloku se odečet sterilní složky neobjeví.

To je pojmenování směru \(\mathcal E_{\mathrm s}\), ne důkaz, že takový stav v přírodě je.
Data ho mohou ukázat.
Model ho nevyrábí.

\section{Co model není}

\begin{remark}
Není dokázáno, že sterilní neutrino existuje.
Není dokázáno, že horizont černé díry takový stav vydává.
Není dokázáno, že Eulerova formule je dodavatel.
Není dokázáno, že nuly funkce \(\xi\) leží na kritické přímce.
Unitární rozšíření není věta o provázání.
\end{remark}

Model popisuje jen toto chování.
Uvnitř se platí.
Zvenčí lze být slyšen bez platby.
Prostor stavů se tím zvětší.
Z nuly, měřeno vnitřním účtem, může být jednotka.
Sterilní směr v ortogonálním součtu je ten, který vnitřní účet nevidí.

\end{document}
